\scanpage{39}\mspage{26}

\diagram{\begin{tikzcd}[ampersand replacement=\&,column sep=large,row sep=large]
\text{(2.4.)}\&
 X_{\deleted{J}\unsure{K|J}}\arrow[r,"\tau^X_{J,K}"] \arrow[d,"u_{\deleted{J}K}"'] \arrow[dr,phantom,"\text{cocart}",pos=.54]
 \& X_{\deleted{K}J}\arrow[d,"\widetilde u_{\deleted{J}K,\deleted{K}J}"]\arrow[dd,bend left=64,"u_{\deleted{K}J}"]\\
\& Y_{\deleted{J}K}\arrow[r,"\tau^Y_{J,K}"]
 \& Y^{\deleted{\unsure{u}}}_{\deleted{J}K,\deleted{K}J}=Y_{\deleted{J}K}\amalg_{X_{\deleted{J}K}}X_{\deleted{K}J}
 \arrow[d,"u_{\deleted{J}K,\deleted{K}J}"]\\
\& \&Y_{\deleted{K}J}
\end{tikzcd}}

\uline{Théorème} 2.1. Soit $M$ \inserted{$\unsure{\mathcal U|\mathscr U}$-}catégorie stable par petites $\unsure{\underrightarrow{\lim}|\underline{\lim}}$, $C\subset\mathrm{Fl}(M)$ satisfaisant les conditions de stabilité (2.1.) (par \unsure{cochgt} de base et par composition transfinie), $I$ un \inserted{petit} ensemble ordonné, $u_\bullet:X_\bullet\longrightarrow Y_\bullet$ une flèche dans $M(I^{\unsure{o|\circ}})$. \deleted{On}

$1^\circ$) \inserted{On suppose $u_\bullet\in C_I$ (cf. (2.2.)).} Soient $\deleted{J}K\subset\deleted{K}J$ deux ouverts de $I$. Alors $u_{\deleted{J}K,\deleted{K}J}$ (défini \unsure{ds|dans} (2.4.)) est dans $C$.

$2^\circ$) \inserted{Supposons $I$ artinien.} Pour qu'on ait $u_\bullet\in C_I$, i.e. pour que $u_{\deleted{J}K,\deleted{K}J}$ soit dans $C$ pour tt couple $(\deleted{J}K\subset\deleted{K}J)$ d'ouverts tels que $\deleted{K}J$ ait un plus grand élément \inserted{i.e. soit de la forme $I_{/i_0}$} (donc aussi, par $1^\circ$, pour \uline{tout} couple d'ouverts $J\subset K$), il (faut et) suffit qu'il en soit ainsi dans le cas particulier où $J=I_{/i_0}$, $J'=I'_{/i_0}$

\[
 \overset{\mathrm{def}}{=} I_{/i_0}-\{i_0\}=\{i\in I\mid i<i_0\}.
\]

% archive sheet number: 38

